\documentclass[10pt]{article}

\usepackage[english]{babel}
\usepackage[letterpaper,top=2cm,bottom=2cm,left=3cm,right=3cm,marginparwidth=1.75cm]{geometry}
\usepackage{helvet}
\renewcommand{\familydefault}{\sfdefault}
\usepackage{setspace}
\setstretch{1.0}
\pagestyle{empty}

\usepackage{amsmath,amssymb,amsfonts}
\usepackage{graphicx}
\usepackage{tabularx}
\usepackage{booktabs}
\usepackage{longtable}
\usepackage{array}
\usepackage{float}
\usepackage{xcolor}
\usepackage[colorlinks=true, allcolors=blue]{hyperref}
\usepackage[labelfont=bf]{caption}

\graphicspath{{figures/111/}}

\renewcommand{\thefigure}{S\arabic{figure}}
\renewcommand{\thetable}{S\arabic{table}}
\renewcommand*{\thesection}{S\arabic{section}}

\newcolumntype{L}[1]{>{\raggedright\arraybackslash}p{#1}}

\title{Supplemental Information for\\
``TomatoPromoterDesigner: a reproducible framework for motif-aware analysis and design of tissue-biased tomato promoters''}

\author{Zhoujie Fan, Haimeng Zhang, Xiaoyu Wang, Hao Meng, Hong Liang,\\
Zheng Wang, Huazhu Fu, Ching-Yu Cheng, Shan Cong*, Xiaohui Yao*}

\date{}

\begin{document}
\maketitle

\section{Supplementary Methods}

\subsection{Software scope and release organization}

TomatoPromoterDesigner is an installable Python command-line framework for tomato promoter sequence processing, motif annotation, tissue-associated prediction, motif-aware candidate design, report generation and figure export. The release contains two complementary layers. The first layer consists of package-native modules that keep the software directly runnable after installation. These include deterministic motif scanning, a deterministic four-tissue baseline predictor, motif-preserving candidate generation, tabular output writing and compact report generation. The second layer consists of documented adapters and retained result tables derived from earlier tomato promoter modeling workflows. These retained resources are used for manuscript-level examples, supplementary figures and provenance-preserved quantitative checks.

This distinction is important for interpreting the manuscript. Package-native commands are intended to demonstrate the current software interface and reproducible workflow. Retained MpraVAE-, DNABERT- and deepseed-derived resources provide historical result context and figure reconstruction support, but they are not presented as newly trained models introduced by the current release. Component boundaries are summarized in Table~\ref{tab:component-boundary}.

\subsection{Data sources and retained result resources}

Repository-distributed data were organized into raw, processed and result-facing resources. Small tabular files required for manuscript reproduction were retained in the repository, whereas large genomes, HDF5 corpora, BLAST databases and heavyweight checkpoints are tracked separately through external manifests. The manuscript-facing retained result pack is generated by:

\begin{verbatim}
make reproduce-legacy
\end{verbatim}

This command executes \texttt{scripts/reproduce\_legacy\_outputs.py} and writes synchronized tables and figures under \texttt{data/results/reproducible\_legacy/}. The generated manifest records the source table and reconstructed output for each retained result. Major data resources used by the supplementary analyses are summarized in Table~\ref{tab:data-resources}. Image file locations are intentionally not listed as a formal table because the figures themselves are part of the submitted supplement; repository image paths are maintained in the project documentation rather than treated as scientific results.

\subsection{Command-line workflow}

The basic workflow begins with promoter sequences in FASTA format and passes them through annotation, prediction, design and report-generation steps:

\begin{verbatim}
pip install -e .

tomato-promoter-designer validate-input \
  --input examples/demo_input.fasta

tomato-promoter-designer annotate \
  --input examples/demo_input.fasta \
  --output outputs/demo_annotate.csv

tomato-promoter-designer predict \
  --input examples/demo_input.fasta \
  --output outputs/demo_predict.csv

tomato-promoter-designer design \
  --input examples/demo_input.fasta \
  --target fruit \
  --candidates 3 \
  --seed 42 \
  --output outputs/demo_design.csv

tomato-promoter-designer report \
  --input outputs/demo_design.csv \
  --output outputs/demo_report.json
\end{verbatim}

The same interface exposes adapter-oriented commands for DNABERT-derived motif processing, MpraVAE-derived four-tissue prediction and design, deepseed-derived scalar prediction, figure export and legacy figure reconstruction (Table~\ref{tab:commands}). The small bundled FASTA example is intended for software testing and documentation. Broader supplementary result summaries use the retained result pack described above.

\subsection{Input validation and sequence normalization}

FASTA records are parsed into sequence identifiers and sequence strings. Empty FASTA files are rejected, empty sequence records are rejected, sequence strings are uppercased, duplicate identifiers are rejected and unsupported symbols outside \texttt{A/C/G/T/N/M} are not accepted. The package-native modules can inspect and score general promoter sequences, including ambiguous \texttt{N} or \texttt{M} symbols where allowed by the validation layer. The migrated MpraVAE-derived routes require canonical 165-bp unambiguous \texttt{A/C/G/T} sequences because this is the input shape expected by the retained model architecture.

\subsection{Package-native motif annotation}

The package-native motif module performs deterministic exact-match scanning using the repository-default motif set \texttt{CAAAA}, \texttt{CTATT}, \texttt{ATTTT}, \texttt{TTAAA} and \texttt{TTTAT}. For every motif occurrence, including overlapping occurrences, the scanner reports the sequence identifier, motif sequence, zero-based start coordinate, zero-based end coordinate and a match score equal to motif length. This module provides a transparent baseline regulatory-interpretation layer and a consistent input to the motif-aware design route.

\subsection{DNABERT-derived motif post-processing}

The DNABERT-derived adapter is exposed as \texttt{annotate-legacy-dnabert}. It processes retained DNABERT sequence and attention-score resources and produces motif-summary, labeled-sequence and metadata outputs. The retained processed motif table is used to summarize the top 20 motif instances in Fig.~\ref{fig:supp-motif}. This route is described as a legacy-derived post-processing workflow because the current Application Note does not introduce a newly trained DNABERT model.

\subsection{Tissue-associated prediction modules}

TomatoPromoterDesigner reports continuous tissue-associated scores for root, stem, leaf and fruit rather than binary tissue labels. The package-native predictor is deterministic and combines GC ratio, AT ratio, poly-A content and default motif counts into four tissue-associated scores. Let \(G_C\) be GC ratio, \(A_T\) be AT ratio, \(P_A\) be normalized \texttt{AAAAA} count and \(c_m\) be the count of motif \(m\). The implemented baseline scores are:

\[
\begin{aligned}
S_{root} &= 4.2G_C + 2.0c_{\mathrm{CTATT}} + 0.5c_{\mathrm{TTAAA}},\\
S_{stem} &= 3.2A_T + 1.8c_{\mathrm{TTTAT}} + 0.6c_{\mathrm{ATTTT}},\\
S_{leaf} &= 2.5G_C + 1.1c_{\mathrm{CAAAA}} + 0.7c_{\mathrm{CTATT}},\\
S_{fruit} &= 3.5A_T + 2.2c_{\mathrm{CAAAA}} + 1.3c_{\mathrm{ATTTT}} + 2.0P_A.
\end{aligned}
\]

The preferred tissue is assigned as the tissue with the highest reported score. When compatible external MpraVAE checkpoints are available and sequences match the expected 165-bp input format, the migrated MpraVAE-derived route can emit four-tissue scores using the retained checkpoint-dependent workflow. Both package-native and adapter-derived routes emit a harmonized output schema (Table~\ref{tab:output-fields}).

\subsection{Motif-aware candidate design}

The package-native design route first annotates motif positions in each input promoter and treats those positions as protected. It then generates a candidate pool by applying seeded substitutions to non-protected positions according to target-tissue nucleotide preferences. The default mutation rate is 0.18 per non-protected position. Candidate pools are scored by the prediction layer and ranked by the requested target-tissue score. For each candidate, the output table records the original sequence, designed sequence, target tissue, candidate rank, four tissue-associated scores, preserved motif summary, design status and number of point differences.

The MpraVAE-derived design adapter provides a retained checkpoint-dependent route when compatible resources are available. It encodes a 165-bp promoter, optimizes the latent representation toward the target-tissue objective, decodes candidate sequences, applies promoter-like quality control and reports ranked unique candidates. This route is included for continuity with earlier tomato promoter modeling workflows and is not required for the package-native baseline workflow.

\subsection{Report generation and figure export}

The \texttt{report} command builds a compact JSON summary from a design table, including row counts, unique designed sequences, target tissues, mutation statistics and best-candidate summaries. The \texttt{figures} and \texttt{legacy-figures} commands export figure summaries from standardized result tables. The supplementary figures were redrawn from retained source tables so that their displayed values correspond to repository-distributed tabular resources rather than isolated image assets.

\subsection{Quantitative reference statistic}

The retained quantitative reference table contains predicted and measured scalar-expression pairs from earlier tomato promoter modeling workflows. Rows with missing, non-numeric or non-finite values in either value column were excluded. No log transformation or additional normalization was applied before the reported statistic. Pearson correlation was calculated on 12,575 valid pairs:

\[
r =
\frac{\sum_i (x_i-\bar{x})(y_i-\bar{y})}
{\sqrt{\sum_i (x_i-\bar{x})^2}\sqrt{\sum_i (y_i-\bar{y})^2}},
\]

where \(x_i\) and \(y_i\) denote retained predicted and measured scalar-expression values, respectively. The resulting value was \(r = 0.852097\), reported as Pearson \(r = 0.852\). This retained table is used as a quantitative traceability resource and figure-reconstruction anchor, not as an independent benchmark of all prediction components in the current release.

\section{Supplementary Results}

\subsection{Retained quantitative reference}

The retained quantitative reference table contains 12,575 valid predicted/measured scalar-expression pairs. Predicted values range from approximately \(-7.04\) to \(7.07\), whereas measured values range from approximately \(-5.99\) to \(8.53\). Predicted and measured values showed a strong positive correlation (Pearson \(r = 0.852097\); Fig.~\ref{fig:supp-quantitative}). This supports the use of the retained table as a quantitative reference for software-level traceability and figure reconstruction.

\subsection{Balanced four-tissue expression subset}

The retained expression heatmap source table contains 100 tomato promoter sequences and four tissue-associated expression values per sequence. The subset is balanced across preferred-tissue classes, with 25 promoters each assigned to fruit, leaf, root and stem according to the dominant retained score. Target-margin values range from 1.624 to 6.034, with a median of 3.6785 and a mean of 3.3995. The heatmap shows visible block structure after grouping by preferred tissue (Fig.~\ref{fig:supp-expression}), supporting the interpretation that the retained table preserves multi-tissue expression structure rather than isolated single-promoter examples.

\subsection{Retained DNABERT-derived motif instances}

The retained DNABERT-derived motif summary reports the 20 most frequent motifs after post-processing. The five most frequent retained motifs are \texttt{TACTT} (1,893 instances), \texttt{CTATT} (1,264 instances), \texttt{ATTTT} (841 instances), \texttt{CAAAA} (722 instances) and \texttt{TTTAT} (384 instances). Nine of the 20 retained motifs are linked to TFBS logo assets in the repository. These results provide a compact motif-level view that complements the package-native exact-match motif scanner (Fig.~\ref{fig:supp-motif}; Table~\ref{tab:motif-top}).

\subsection{4-mer similarity between original and generated promoters}

The retained 4-mer comparison table contains 256 4-mers for each of three regions: the full promoter sequence, the proximal 20-bp segment and the distal 20-bp segment. Correlation between original and generated 4-mer frequencies is strongest in the proximal 20-bp region (\(r = 0.835\)), followed by the full promoter region (\(r = 0.767\)) and the distal 20-bp region (\(r = 0.585\)). These results indicate that retained generated sequences preserve substantial sequence-composition structure, with regional differences visible in the reconstructed panels (Fig.~\ref{fig:supp-kmer}; Table~\ref{tab:kmer}).

\subsection{Retained fruit-targeted design candidates}

The retained design summary table contains 199 fruit-targeted promoter candidates. Predicted fruit-associated scores range from 1.7999 to 18.4738, with a median of 6.4398 and a mean of 6.7601. Fruit-associated bias margins range from 0.8936 to 2.9709, with a median of 1.8274 and a mean of 1.8522. The number of point differences between original and generated sequences ranges from 82 to 131, with median and mean values both equal to 105. These distributions provide a broader retained design-output view than the two-promoter example described in the main text (Fig.~\ref{fig:supp-design}; Table~\ref{tab:design-summary}).

\section{Supplementary Tables}

\begin{longtable}{L{0.22\textwidth}L{0.28\textwidth}L{0.39\textwidth}}
\caption{Repository data resources used by the manuscript and supplementary analyses.}\label{tab:data-resources}\\
\toprule
Resource & Role in this work & Use in manuscript or supplement \\
\midrule
\endfirsthead
\toprule
Resource & Role in this work & Use in manuscript or supplement \\
\midrule
\endhead
Bundled demo FASTA & Small package-native input file containing two tomato promoter examples & Installation test, quickstart workflow and main-text software demonstration. \\
Tomato promoter training table & Retained promoter sequence and four-tissue expression table & Source for the balanced four-tissue heatmap subset in Fig.~\ref{fig:supp-expression}. \\
Generated prediction reference table & Retained predicted/measured scalar-expression pairs & Quantitative reference in Fig.~\ref{fig:supp-quantitative} and Table~\ref{tab:quant-reference}. \\
Retained DNABERT motif summary & Processed motif table from the DNABERT-derived post-processing route & Top motif summary in Fig.~\ref{fig:supp-motif} and Table~\ref{tab:motif-top}. \\
Generated promoter table & Retained generated promoter candidates from earlier design workflows & 4-mer comparison and design-candidate summary in Figs.~\ref{fig:supp-kmer} and \ref{fig:supp-design}. \\
Reproducible retained result pack & Regenerated tables and figures produced by \texttt{make reproduce-legacy} & Source of the final supplementary figure data and summary statistics. \\
External artifact manifest & List of heavyweight resources that are not bundled by default & Documents checkpoints, large raw resources and files requiring separate release handling. \\
\bottomrule
\end{longtable}

\begin{longtable}{L{0.23\textwidth}L{0.26\textwidth}L{0.40\textwidth}}
\caption{Command-line modules referenced by the manuscript and supplementary materials.}\label{tab:commands}\\
\toprule
Command & Main input & Main output \\
\midrule
\endfirsthead
\toprule
Command & Main input & Main output \\
\midrule
\endhead
\texttt{validate-input} & FASTA promoter sequences & Terminal validation summary and error reporting for invalid inputs. \\
\texttt{annotate} & FASTA promoter sequences & CSV motif-hit table with sequence identifier, motif, start, end and score. \\
\texttt{annotate-legacy-dnabert} & Retained DNABERT sequence and attention resources & Processed motif-summary table, labeled sequence table and metadata output. \\
\texttt{predict} & FASTA promoter sequences & CSV table with root, stem, leaf and fruit scores plus preferred tissue. \\
\texttt{predict-legacy-mpravae} & FASTA plus optional compatible checkpoint & Four-tissue retained checkpoint-dependent prediction table when available. \\
\texttt{predict-legacy-deepseed} & FASTA plus optional compatible checkpoint & Scalar-expression retained prediction table when available. \\
\texttt{design} & FASTA, target tissue and candidate number & Candidate table with designed sequences, scores, point differences and QC flags. \\
\texttt{design-legacy-mpravae} & FASTA plus optional compatible checkpoint & Retained latent-space design route when compatible resources are available. \\
\texttt{report} & Design output table & Compact JSON summary of candidates, mutation counts and best candidates. \\
\texttt{figures} & Standardized result table & Lightweight figure summaries from prediction, design or motif outputs. \\
\texttt{legacy-figures} & Retained historical result tables & Reconstructed retained figure bundle and manifest. \\
\texttt{make reproduce-legacy} & Repository-distributed retained tables & Regenerated retained result pack used for supplementary reporting. \\
\bottomrule
\end{longtable}

\begin{table}[H]
\centering
\caption{Package-native baseline predictor features and weights.}\label{tab:baseline-predictor}
\begin{tabular}{L{0.16\textwidth}L{0.67\textwidth}}
\toprule
Tissue score & Implemented deterministic scoring rule \\
\midrule
Root & \(4.2G_C + 2.0c_{\mathrm{CTATT}} + 0.5c_{\mathrm{TTAAA}}\) \\
Stem & \(3.2A_T + 1.8c_{\mathrm{TTTAT}} + 0.6c_{\mathrm{ATTTT}}\) \\
Leaf & \(2.5G_C + 1.1c_{\mathrm{CAAAA}} + 0.7c_{\mathrm{CTATT}}\) \\
Fruit & \(3.5A_T + 2.2c_{\mathrm{CAAAA}} + 1.3c_{\mathrm{ATTTT}} + 2.0P_A\) \\
\bottomrule
\end{tabular}
\end{table}

\begin{longtable}{L{0.22\textwidth}L{0.66\textwidth}}
\caption{Manuscript-facing output fields.}\label{tab:output-fields}\\
\toprule
Field & Meaning \\
\midrule
\endfirsthead
\toprule
Field & Meaning \\
\midrule
\endhead
\texttt{sequence\_id} & Identifier assigned to an input promoter sequence. \\
\texttt{sequence} & Normalized promoter sequence used for prediction or annotation. \\
\texttt{motif} & Detected motif sequence in package-native annotation or retained motif-summary outputs. \\
\texttt{start}, \texttt{end} & Zero-based motif coordinates reported by the package-native exact-match scanner. \\
\texttt{score} & Motif match score; equal to motif length for package-native exact matching. \\
\texttt{expr\_root}, \texttt{expr\_stem}, \texttt{expr\_leaf}, \texttt{expr\_fruit} & Continuous tissue-associated prediction scores emitted on the common four-tissue schema. \\
\texttt{preferred\_tissue} & Tissue with the largest reported score for a promoter. \\
\texttt{target\_tissue} & User-specified objective for candidate design. \\
\texttt{candidate\_rank} & Rank of a designed candidate within the candidate set for one input promoter. \\
\texttt{original\_sequence}, \texttt{designed\_sequence} & Input promoter and generated candidate sequence. \\
\texttt{preserved\_motifs} & Summary of motifs protected during package-native motif-aware design. \\
\texttt{num\_mutations} or \texttt{num\_point\_differences} & Number of point differences between original and generated promoter sequences. \\
\texttt{passes\_qc} & Route-specific quality-control indicator when available. \\
\bottomrule
\end{longtable}

\begin{table}[H]
\centering
\caption{Balanced retained four-tissue heatmap subset.}\label{tab:heatmap-subset}
\begin{tabular}{lcc}
\toprule
Preferred tissue & Number of promoters & Selection description \\
\midrule
Fruit & 25 & Dominant retained fruit score \\
Leaf & 25 & Dominant retained leaf score \\
Root & 25 & Dominant retained root score \\
Stem & 25 & Dominant retained stem score \\
\midrule
Total & 100 & Balanced retained subset used in Fig.~\ref{fig:supp-expression} \\
\bottomrule
\end{tabular}
\end{table}

\begin{table}[H]
\centering
\caption{Top retained DNABERT-derived motif instances.}\label{tab:motif-top}
\begin{tabular}{lcc}
\toprule
Motif & Retained instances & TFBS logo linked \\
\midrule
\texttt{TACTT} & 1,893 & No \\
\texttt{CTATT} & 1,264 & Yes \\
\texttt{ATTTT} & 841 & Yes \\
\texttt{CAAAA} & 722 & Yes \\
\texttt{TTTAT} & 384 & Yes \\
\texttt{TCTTA} & 182 & No \\
\texttt{ACTAG} & 145 & No \\
\texttt{TTAAA} & 120 & Yes \\
\texttt{GCTTT} & 116 & No \\
\texttt{TCCGA} & 111 & No \\
\bottomrule
\end{tabular}
\end{table}

\begin{table}[H]
\centering
\caption{Retained 4-mer frequency comparison by sequence region.}\label{tab:kmer}
\begin{tabular}{lcc}
\toprule
Region & Number of 4-mers & Pearson correlation between original and generated frequencies \\
\midrule
Full promoter sequence & 256 & 0.767 \\
Proximal 20 bp & 256 & 0.835 \\
Distal 20 bp & 256 & 0.585 \\
\bottomrule
\end{tabular}
\end{table}

\begin{table}[H]
\centering
\caption{Retained fruit-targeted design-candidate summary.}\label{tab:design-summary}
\begin{tabular}{lcccc}
\toprule
Metric & Minimum & Median & Mean & Maximum \\
\midrule
Predicted fruit-associated score & 1.7999 & 6.4398 & 6.7601 & 18.4738 \\
Fruit-associated bias margin & 0.8936 & 1.8274 & 1.8522 & 2.9709 \\
Number of point differences & 82 & 105 & 105 & 131 \\
\bottomrule
\end{tabular}
\end{table}

\begin{table}[H]
\centering
\caption{Retained quantitative reference statistics.}\label{tab:quant-reference}
\begin{tabular}{lc}
\toprule
Statistic & Value \\
\midrule
Valid predicted/measured pairs & 12,575 \\
Pearson correlation & 0.852097 \\
Reported Pearson correlation & 0.852 \\
Predicted value range & \(-7.0409\) to \(7.0720\) \\
Measured value range & \(-5.9943\) to \(8.5316\) \\
Transformation before correlation & None \\
Interpretation & Retained quantitative reference; not a benchmark of all prediction backends \\
\bottomrule
\end{tabular}
\end{table}

\begin{longtable}{L{0.22\textwidth}L{0.28\textwidth}L{0.39\textwidth}}
\caption{Component availability and interpretation boundary.}\label{tab:component-boundary}\\
\toprule
Component & Availability in current release & Interpretation in this manuscript \\
\midrule
\endfirsthead
\toprule
Component & Availability in current release & Interpretation in this manuscript \\
\midrule
\endhead
Package-native motif scanner & Included in source code & Deterministic exact-match motif annotation and motif-protection support. \\
Package-native predictor & Included in source code & Deterministic baseline scorer for runnable four-tissue outputs. \\
Package-native designer & Included in source code & Motif-aware candidate generation using protected motif positions and seeded substitutions. \\
MpraVAE-derived prediction/design adapter & Source adapter included; compatible checkpoints may be external & Retained checkpoint-dependent route for 165-bp tomato promoter prediction and design. \\
DNABERT-derived motif adapter & Source adapter and selected support resources included & Retained attention-to-motif post-processing route when support files are supplied. \\
deepseed-derived scalar predictor & Source adapter included; compatible checkpoint may be external & Retained scalar-expression prediction route for historical compatibility. \\
Retained quantitative reference table & Included as a repository-distributed retained table & Quantitative traceability resource and figure-reconstruction anchor. \\
Large model artifacts & Tracked through external manifests rather than bundled by default & Optional resources requiring separate availability and redistribution handling. \\
\bottomrule
\end{longtable}

\section{Supplementary Figures}

\begin{figure}[H]
\centering
\includegraphics[width=0.82\textwidth]{Fig1D_quantitative_reference.png}
\caption{Retained quantitative reference scatter plot. Each point represents one predicted/measured scalar-expression pair retained from earlier tomato promoter modeling workflows (\(n = 12{,}575\)). The dashed line indicates the identity relationship. The retained reference set showed a strong positive correlation between predicted and measured scalar-expression values (Pearson \(r = 0.852\)). This panel is presented as a retained quantitative reference and not as an independent benchmark of all prediction components in the current software release.}
\label{fig:supp-quantitative}
\end{figure}

\begin{figure}[H]
\centering
\includegraphics[width=0.72\textwidth]{FigS2_expression_heatmap_final.png}
\caption{Retained tissue-associated expression profiles for a balanced 100-promoter subset. Columns correspond to root, stem, leaf and fruit expression values preserved in the retained source table. Rows are grouped by \texttt{preferred\_tissue}, with 25 promoters represented in each class. The side color strip marks the preferred-tissue label assigned from the dominant retained score. This panel visualizes the structure of the bundled four-tissue expression outputs rather than a standalone benchmark.}
\label{fig:supp-expression}
\end{figure}

\begin{figure}[H]
\centering
\includegraphics[width=0.82\textwidth]{FigS3_dnabert_motif_summary.png}
\caption{Retained DNABERT-derived motif summary. Bars show the number of retained instances for the 20 most frequent sequence motifs preserved after DNABERT-derived post-processing. The panel summarizes retained motif patterns used to support inspection of the motif-analysis component described in the main text.}
\label{fig:supp-motif}
\end{figure}

\begin{figure}[H]
\centering
\includegraphics[width=\textwidth]{FigS4_kmer_distribution.png}
\caption{Retained 4-mer frequency comparisons between original and generated tomato promoter sequences. Each point represents one 4-mer. Panels summarize the full promoter sequence, the proximal 20-bp region and the distal 20-bp region, with dashed lines indicating identity relationships. Pearson correlation coefficients quantify the similarity between retained original and generated 4-mer frequency profiles in each region.}
\label{fig:supp-kmer}
\end{figure}

\begin{figure}[H]
\centering
\includegraphics[width=\textwidth]{figS5_design_candidate_summary.png}
\caption{Retained fruit-targeted design-candidate summary. (A) Distribution of predicted fruit-associated scores across retained candidate promoters. (B) Distribution of fruit-associated bias margins, defined as the difference between the retained fruit-associated score and the strongest non-fruit tissue-associated score. (C) Distribution of the number of point differences between original and generated promoter sequences. Dashed lines indicate median values for each summary distribution.}
\label{fig:supp-design}
\end{figure}

\end{document}
